\section{Discussion}
FSR is a robust metric because it is mainly set by round-trip length and group index. Once the ring radius and waveguide dispersion are fixed, FSR changes smoothly and is less sensitive to small loss changes than linewidth-based metrics. In contrast, \(Q\) and finesse depend strongly on propagation loss, bend loss, coupling strength, spectral sampling resolution, mesh resolution, monitor placement, and numerical convergence. The lab notes explicitly caution that \(Q\) factor and finesse are sensitive to loss and minor numerical or fabrication variations \cite{labmanual}.

INTERCONNECT, FDE, and varFDTD represent different abstraction levels. INTERCONNECT is fast and circuit-oriented, making it useful for transfer functions and parameter sweeps. FDE resolves the waveguide cross-section and provides modal quantities such as \(n_{\mathrm{eff}}\), \(n_g\), and coupled-mode index splitting. varFDTD verifies the final geometry with a broadband propagation simulation under an effective-index approximation. Agreement among the three levels gives confidence; disagreement usually points to missing loss, dispersion, insufficient meshing, or an inconsistent coupling model.

For sensing, the resonator converts a measurand into a spectral shift. Analyte binding, refractive-index change, temperature, or strain modifies \(n_{\mathrm{eff}}\) or the optical path length. Practical interrogation can use wavelength scanning for direct spectra, resonance tracking for high resolution, or fixed-wavelength intensity readout for compact instrumentation. The fixed-wavelength approach is efficient but more vulnerable to laser power fluctuation unless differential through/drop readout or amplitude-comparison techniques are used.

Fabrication tolerance is critical. A small error in waveguide width or height changes dispersion and therefore \(n_g\), shifting the FSR and resonance locations. Gap variation changes \(|k|^2\), linewidth, extinction ratio, and drop efficiency. Sidewall roughness and bend loss reduce \(Q\), while minimum lithographic gap rules may constrain the strongest practical coupling. A complete design should therefore include mesh refinement, a gap and width tolerance sweep, and experimental calibration.

The available targets also contain a useful consistency check. The target \(Q=1800\) at \SI{1550}{\nano\metre} implies \(FWHM=\lambda/Q=\SI{0.861}{\nano\metre}\). The target finesse of 26 with \(FSR=\SI{26.2}{\nano\metre}\) implies \(FWHM=FSR/\mathcal{F}=\SI{1.008}{\nano\metre}\). These are close in scale but not identical, so the final design should specify whether the \(Q\) target or finesse target is prioritised after real spectra are exported.
